\chapter{Étude préalable du sujet}

\section*{Introduction}
\addcontentsline{toc}{section}{Introduction}

Dans ce chapitre introductif, nous définissons le cadre global de notre projet.
Nous présentons d'abord l'organisme d'accueil, la Banque de Tunisie et des Emirats,
au sein de laquelle ce stage a été réalisé. Nous décrivons ensuite le processus de
revue de sécurité tel qu'il existe aujourd'hui à la BTE, avant d'en pointer les
limites et de présenter la solution proposée. Nous exposons enfin la méthodologie
adoptée, du choix du framework Agile au déroulement effectif du projet sur cinq mois.

% ─────────────────────────────────────────────────────────────────────────────
\section{Présentation de l'organisme d'accueil}

    \subsection{Présentation générale}

        La Banque de Tunisie et des Emirats (BTE) \cite{BTE} a été fondée le 10 août 1982 à la suite d'une
        convention signée entre le gouvernement tunisien et l'Abu Dhabi Investment
        Authority (ADIA). Conçue à l'origine comme banque d'investissement, elle est
        passée en 2004 au statut de banque universelle et propose désormais une gamme
        complète de produits et services financiers aux particuliers, aux professionnels
        et aux grandes entreprises. Son actionnariat binational lui donne une couverture
        forte sur les opérations de commerce international, et sa plateforme digitale
        \textbf{NEO BTE} traduit son orientation actuelle vers la modernisation des
        canaux client. Le siège social, situé au Centre Urbain Nord de Tunis, accueille
        un réseau d'agences en cours d'extension.

        Le tableau \ref{tab:bte_fiche} présente la fiche d'identité de l'établissement.

        \begin{table}[H]
        \centering
        \begin{tabular}{|m{4.5cm}|m{8.5cm}|}
        \hline
        \textbf{Dénomination} & Banque de Tunisie et des Emirats (BTE) \\
        \hline
        \textbf{Date de création} & 10 août 1982 \\
        \hline
        \textbf{Forme juridique} & Société anonyme cotée en bourse \\
        \hline
        \textbf{Actionnariat} & Partenariat entre l'État Tunisien et l'Abu Dhabi
                                Investment Authority (ADIA) \\
        \hline
        \textbf{Siège social} & Boulevard Béji Caïed Essebsi, Centre Urbain Nord,
                                1082 Tunis \\
        \hline
        \textbf{Secteur d'activité} & Banque universelle (Particuliers, Professionnels,
                                      Entreprises) \\
        \hline
        \textbf{Produits clés} & Crédits, épargne, e-banking (NEO BTE), services
                                 internationaux, leasing \\
        \hline
        \end{tabular}
        \caption{Fiche d'identité de la Banque de Tunisie et des Emirats}
        \label{tab:bte_fiche}
        \end{table}

        Dans le cadre de sa transformation numérique, la BTE a mis en place une Direction
        des Systèmes d'Information (DSI). Celle-ci est chargée d'assurer la disponibilité,
        la sécurité et l'évolution continue de ses infrastructures informatiques. Notre
        stage s'est déroulé au sein du département Informatique et Digital de cette
        direction, responsable de la modernisation des systèmes d'information et de
        la plateforme NEO BTE. La mission qui nous a été confiée est de renforcer les
        pratiques DevSecOps de l'organisation.

    \subsection{Secteur d'activité}

        L'activité de la BTE se déploie sur trois domaines qui couvrent l'ensemble des
        besoins financiers de ses clients. Le premier est la banque de détail :
        comptes courants, épargne, crédits à la consommation, cartes bancaires et packs
        dédiés (Packs Professionnels, Médecins, etc.). Le second est le financement des
        entreprises, à travers des crédits d'investissement, des lignes de crédit, des
        opérations de leasing et des services de commerce international (lettres de
        crédit, remises documentaires, transferts SWIFT) ; sur ce dernier point, la BTE
        a construit une expertise reconnue sur le marché tunisien grâce à son
        actionnariat binational. Le troisième domaine est la transformation digitale
        autour de la plateforme \textbf{NEO BTE}, qui rénove progressivement les
        interfaces client et les systèmes d'information de la banque. C'est cette
        dynamique de modernisation qui a rendu pertinent notre stage : renforcer la
        sécurité des développements logiciels en s'appuyant sur l'intelligence
        artificielle locale.

    \subsection{Organigramme}

        \begin{figure}[H]
            \centering
            \frame{\includegraphics[width=0.85\columnwidth]{organigramme_dsi}}
            \caption{Organigramme de la Direction des Systèmes d'Information — BTE}
            \label{fig:organigramme}
        \end{figure}

        La figure \ref{fig:organigramme} présente l'organigramme de la Direction des
        Systèmes d'Information, organisée en plusieurs pôles fonctionnels : développement
        applicatif, infrastructure et réseaux, sécurité informatique et support utilisateurs.
        Notre projet de stage s'est déroulé au sein du pôle Sécurité Informatique, en
        collaboration avec les équipes de développement et d'infrastructure afin de mettre
        en place une solution DevSecOps transversale à l'ensemble de la DSI.

% ─────────────────────────────────────────────────────────────────────────────
\section{Présentation du projet}

    Le projet vise à automatiser la revue de sécurité du code à la BTE. Il associe
    pour cela de l'intelligence artificielle locale et des outils d'analyse statique. L'agent
    prend en charge l'ensemble du cycle, du déclenchement des scanners à la publication
    des commentaires de correction sur la \textit{Pull Request}, sans intervention
    humaine entre les deux. L'objectif est double : raccourcir un délai de revue qui
    dépasse aujourd'hui quarante-huit heures, et garantir qu'aucune \textit{Pull
    Request} ne soit fusionnée sans avoir été passée au crible des cinq scanners.

    \subsection{Étude de l'existant}

        Le processus de mise en production des développements logiciels à la BTE
        s'articule autour de la plateforme GitHub, qui sert à la fois de système de
        gestion de versions et d'espace de revue de code. Le cycle commence par la
        création d'une branche dédiée à partir de la branche principale du dépôt.
        Le développeur y pousse ses commits au fil de l'implémentation, puis ouvre
        une \textit{Pull Request} vers la branche cible une fois la fonctionnalité
        prête pour relecture.

        À l'ouverture de la \textit{Pull Request}, GitHub notifie l'équipe par le
        mécanisme des \textit{code owners} configuré sur les répertoires concernés.
        Selon la nature des fichiers modifiés, un développeur senior ou un référent
        sécurité prend la relecture en charge dans l'interface GitHub : examen visuel
        du diff, lecture des fichiers touchés, recoupements éventuels avec le reste
        du code base. Les remarques sont déposées sous forme de commentaires inline
        rattachés aux lignes concernées, et le dialogue avec le développeur se
        poursuit dans le fil de la \textit{Pull Request}.

        Le développeur prend en compte les retours, pousse des commits correctifs
        sur la même branche, et la \textit{Pull Request} se met automatiquement à
        jour. Une fois le code jugé prêt par le relecteur, la fusion vers la branche
        principale est déclenchée manuellement. La revue repose intégralement sur
        l'examen humain ; aucune analyse statique automatisée, aucun scanner de
        vulnérabilités et aucun \textit{gate} de pipeline ne sont greffés sur ce
        cycle. Le résultat de chaque revue est conservé dans l'historique GitHub
        de la \textit{Pull Request} concernée, accessible \textit{a posteriori} par
        l'équipe.

    \subsection{Critique de l'existant}

        L'analyse du processus existant fait apparaître plusieurs limites qui justifient
        une automatisation. Les délais de revue dépassent fréquemment quarante-huit heures,
        ce qui ralentit la livraison des fonctionnalités et laisse des vulnérabilités
        exposées plus longtemps que nécessaire. La qualité de la revue elle-même dépend
        de la disponibilité et de l'expertise du relecteur du jour : selon les
        \textit{Pull Requests}, la couverture de sécurité varie sensiblement d'un cas à
        l'autre. L'absence de \textit{gate} automatisé dans le pipeline CI/CD est sans
        doute la limite la plus directe : une \textit{Pull Request} contenant une
        vulnérabilité critique peut aujourd'hui être fusionnée sans aucun blocage
        technique. Et la gestion entièrement manuelle du processus laisse régulièrement
        passer des catégories de risques pour lesquelles l'œil humain n'est pas le bon
        outil, en particulier les secrets accidentellement commités et les dépendances
        vulnérables introduites en transitif.

        \begin{figure}[H]
            \centering
            \frame{\includegraphics[width=\columnwidth,height=0.80\textheight,keepaspectratio]{asis_workflow}}
            \caption{Processus \textit{AS-IS} — revue manuelle de \textit{Pull
                     Request} telle qu'elle existe aujourd'hui à la BTE.}
            \label{fig:asis_workflow}
        \end{figure}



    \subsection{Solution proposée}

        \begin{figure}[H]
            \centering
            \frame{\includegraphics[width=\textwidth]{solution_overview}}
            \caption{Vue d'ensemble de la solution BTE Security AI Agent}
            \label{fig:solution_overview}
        \end{figure}

        Face à ces défis, nous proposons le \textbf{BTE Security AI Agent}, un moteur
        d'automatisation DevSecOps piloté par l'intelligence artificielle locale, dont la
        vue d'ensemble est présentée dans la figure \ref{fig:solution_overview}. La solution
        repose sur l'interception des événements GitHub via un webhook sécurisé HMAC-SHA256,
        l'orchestration d'un workflow LangGraph à neuf nœuds, et l'utilisation de deux
        modèles de langage locaux exécutés entièrement via Ollama sur le VPS de la BTE.
        Cinq outils d'analyse s'exécutent en parallèle selon la classification de la
        \textit{Pull Request} : Trivy pour les CVE, Gitleaks pour la détection de secrets,
        Semgrep pour les failles OWASP, Checkov pour les misconfigurations d'infrastructure,
        et OSV-Scanner pour les dépendances vulnérables. Le résultat est publié directement
        sur la \textit{Pull Request} GitHub en environ 15 à 25 minutes selon la complexité
        de la PR. Cette publication comprend une revue de sécurité complète avec score de
        risque, verdict et commentaires inline accompagnés de suggestions de correction.
        La totalité de l'inférence s'exécute en local. Aucun code source ne quitte
        l'infrastructure de la BTE.

        \begin{figure}[H]
            \centering
            \frame{\includegraphics[width=\columnwidth,height=0.85\textheight,keepaspectratio]{tobe_workflow}}
            \caption{Processus \textit{TO-BE} — revue automatisée par le BTE
                     Security AI Agent, déclenchée à chaque \textit{Pull Request}.}
            \label{fig:tobe_workflow}
        \end{figure}
% ─────────────────────────────────────────────────────────────────────────────
\section{Approche méthodologique}

    Le choix d'une méthodologie de travail conditionne directement la capacité d'un
    projet à tenir ses délais sans pour autant ignorer les imprévus qui apparaissent dès
    qu'on déploie en production. Dans notre cas, nous voulions une structure suffisamment
    cadrée pour rester lisible côté encadrement, et suffisamment souple pour absorber
    les corrections d'urgence sans casser la planification.

    \subsection{Généralités}

        Le cycle de vie d'un projet logiciel couvre l'analyse de l'existant, la
        spécification des besoins, la conception, la réalisation, les tests,
        l'exploitation et la maintenance. Plusieurs familles de méthodes encadrent ce
        cycle : le cycle en V, qui découpe le projet en phases successives ; le découpage
        en activités de type WBS ; et les méthodes Agiles, qui privilégient les
        itérations courtes et l'ajustement continu. Les méthodes classiques restent
        vulnérables à la mauvaise interprétation des besoins, aux changements
        fonctionnels en cours de route et aux dépassements de délais ; les chiffres
        régulièrement cités dans la littérature font état d'environ un quart de projets
        réussis sous ces approches \cite{Agile}, ce qui rend l'option Agile plus
        pertinente dans le contexte de notre stage.

    \subsection{Comparaison des frameworks Agile}

        Le choix de l'approche Agile étant acté, encore faut-il sélectionner le
        framework le mieux adapté à un projet d'ingénieur en solo, déployé en
        environnement de production réel, sur une durée de cinq mois avec des
        jalons d'encadrement fixés à l'avance. La famille Agile compte plusieurs
        cadres dont la philosophie diverge sensiblement. Le tableau
        \ref{tab:agile_compare} en propose une synthèse comparée sur les critères
        qui pèsent le plus dans notre contexte.

        \begin{table}[H]
        \centering
        \footnotesize
        \begin{tabular}{|m{2cm}|m{2.6cm}|m{2.6cm}|m{2.6cm}|m{2.8cm}|}
        \hline
        \textbf{Critère} & \textbf{Scrum} & \textbf{Kanban} & \textbf{XP} & \textbf{Scrumban} \\
        \hline
        Cadence
        & Sprints fixes (1 à 4 semaines)
        & Flux continu, pas de sprint
        & Itérations très courtes (1 semaine)
        & Sprints + flux Kanban interne \\
        \hline
        Planification
        & \textit{Sprint planning} avec engagement par lot
        & \textit{Pull} à la demande, aucun engagement
        & \textit{Planning game} basé sur la vélocité
        & Planning léger + \textit{pull} intra-sprint \\
        \hline
        Limites WIP
        & Implicite (taille du sprint)
        & Explicite par colonne
        & Implicite (binôme par tâche)
        & Explicite par colonne \\
        \hline
        Cérémonies
        & Daily, planning, démo, rétrospective
        & Aucune obligatoire
        & \textit{Pair programming}, TDD, refactoring constant
        & Cérémonies Scrum allégées \\
        \hline
        Rôles
        & PO, Scrum Master, équipe
        & Aucun rôle imposé
        & Coach, customer, programmer
        & Hérités de Scrum, simplifiés \\
        \hline
        Adapté solo
        & Difficile (rôles redondants)
        & Oui
        & Inadapté (\textit{pair} central)
        & Oui \\
        \hline
        Gestion imprévus
        & Faible — change le sprint en cours
        & Forte — \textit{pull} immédiat
        & Moyenne
        & Forte — classe \textit{Expedite} et WIP \\
        \hline
        Visibilité jalons
        & Forte (dates de sprint fixes)
        & Faible (flux continu)
        & Forte (\textit{releases} fréquentes)
        & Forte (sprints conservés) \\
        \hline
        \end{tabular}
        \caption{Comparaison synthétique des frameworks Agile évalués}
        \label{tab:agile_compare}
        \end{table}

        Scrum, dans sa forme canonique, suppose une équipe pluridisciplinaire avec
        des rôles distincts (Product Owner, Scrum Master, équipe de développement) ;
        ce découpage perd son sens lorsque l'équipe se réduit à une seule personne,
        et le sprint engagé devient rigide dès qu'un incident de production
        s'impose en milieu de quinzaine. \textit{eXtreme Programming} pousse plus
        loin encore les pratiques d'ingénierie (\textit{pair programming}, TDD strict,
        refactoring continu), dont l'essentiel est inopérant en solo.
        Kanban offre la souplesse maximale du flux continu, mais sacrifie la
        cadence et les jalons fixes utiles à l'encadrement académique d'un projet
        de fin d'études.

        Le \textbf{Scrumban} \cite{Ladas2009} a été retenu parce qu'il combine ce
        que Scrum et Kanban font de mieux : la cadence prévisible et les jalons de
        Scrum, couplés à la discipline de flux et aux limites WIP de Kanban. Ce
        compromis répond précisément à la double contrainte du stage : tenir un
        planning lisible pour l'encadrement aux jalons fixés à l'avance, sans pour
        autant figer la gestion quotidienne lorsqu'un incident de production
        s'invite dans le sprint en cours.

    \subsection{Adaptation de Scrumban au stage}

        Formalisée à la fin des années 2000 dans les équipes lean, la méthode
        Scrumban combine la cadence prévisible de Scrum à la discipline de flux
        propre à Kanban. Au lieu de pousser un lot de tâches au début du sprint
        et d'espérer qu'il tienne jusqu'au bout, l'équipe « tire » la prochaine
        carte dès qu'une place se libère, dans la limite imposée par le
        \textit{Work In Progress} (WIP). Le cadre obtenu garde la lisibilité du
        planning bisemestriel, utile pour l'encadrement et les jalons, tout en
        absorbant les imprévus de production sans avoir à replanifier le sprint
        en cours.

        Cette mécanique s'est imposée naturellement pour un stage d'ingénieur en
        solo, dans un environnement où les incidents arrivent rarement au moment
        où on les attend. Le rythme reste celui de Scrum \cite{Scrum} : sprints
        de deux semaines, \textit{sprint planning} en début de quinzaine, démo
        et rétrospective à la fin. Les cérémonies prévues pour les équipes
        (daily standup formel, planning poker, Scrum Master dédié) ont été
        simplifiées. La structure fondamentale a toutefois été conservée :
        backlog priorisé, sprints à durée fixe, jalons de livraison et
        définition de terminé. La figure \ref{fig:agile_scrum} illustre ce cycle.

        \begin{figure}[H]
            \centering
            \frame{\includegraphics[width=0.82\columnwidth]{agile_scrum_cycle}}
            \caption{Cycle Scrumban adapté — sprints de deux semaines avec jalons et flux Kanban interne. \\
                     \footnotesize\textit{Source : Wikimedia Commons, CC0.}}
            \label{fig:agile_scrum}
        \end{figure}

        La gestion au jour le jour suit en revanche la logique Kanban \cite{Kanban}. Une
        seule tâche est en cours à tout moment (WIP=1), entièrement déployée et validée en
        production avant que la suivante ne démarre. La figure \ref{fig:kanban_board}
        illustre la forme générale d'un tableau de ce type, organisé en colonnes
        successives avec des plafonds par colonne.

        \begin{figure}[H]
            \centering
            \frame{\includegraphics[width=0.88\columnwidth]{kanban_board}}
            \caption{Tableau Kanban — visualisation du flux avec limites WIP par colonne. \\
                     \footnotesize\textit{Source : Jennifer Falco, Wikimedia Commons, CC BY 4.0.}}
            \label{fig:kanban_board}
        \end{figure}

        La règle WIP=1 a un effet pratique direct : elle empêche le multitâche
        silencieux, cette dérive bien connue où plusieurs cartes traînent à 80\% sans
        qu'aucune ne franchisse la ligne d'arrivée. Elle réserve aussi une voie nette
        pour les incidents de production. Quand l'urgence disque a saturé le VPS le
        20 avril, quand AlertManager s'est révélé silencieusement cassé depuis le
        déploiement, quand VictoriaMetrics est tombée pour neuf jours, chacune de ces
        découvertes est devenue la carte prioritaire du moment. Le sprint en cours
        absorbait l'imprévu ; le sprint suivant intégrait l'amélioration durable qui
        en découlait. La structure globale du projet n'a pas eu à être revue.

    \subsection{Structure du tableau et choix des colonnes}

        Le tableau Scrumban du projet a été configuré dans Jira en partant d'un projet
        Scrum classique, auquel ont été ajoutés la fonctionnalité \textit{Kanban backlog}
        et des plafonds WIP explicites par colonne. Six colonnes composent le flux. Le
        tableau \ref{tab:scrumban_columns} récapitule leur rôle et la limite associée.

        \begin{table}[H]
        \centering
        \begin{tabular}{|m{2.2cm}|m{2.6cm}|c|m{6.4cm}|}
        \hline
        \textbf{Colonne} & \textbf{Statut Jira} & \textbf{WIP} & \textbf{Rôle} \\
        \hline
        Backlog & \texttt{To Do} & — &
            Réservoir d'idées, de dette technique et de propositions, volontairement non limité \\
        \hline
        Ready & \texttt{Ready} & 5 &
            Sprint backlog issu du \textit{sprint planning}, plafond indicatif contre le sur-engagement \\
        \hline
        In Progress & \texttt{In Progress} & \textbf{1} &
            Carte en cours de développement actif --- limite WIP=1 structurante, empêche le multitâche silencieux \\
        \hline
        In Review & \texttt{In Review} & 2 &
            Revue de code ou validation fonctionnelle avant fusion ; le plafond de 2 évite l'interruption du développement \\
        \hline
        Testing & \texttt{Testing} & 2 &
            Vérification en production sur événement réel (webhook PR, alerte Prometheus, requête chat) \\
        \hline
        Done & \texttt{Done} & — &
            Carte ayant satisfait les cinq conditions de la définition de terminé ; sinon retour en \textit{In Progress} ou \textit{Testing} \\
        \hline
        \end{tabular}
        \caption{Colonnes du tableau Scrumban et limites WIP appliquées}
        \label{tab:scrumban_columns}
        \end{table}

        Une étiquette \textbf{Expedite}, portée par les cartes d'incident de production,
        autorise une exception explicite à la limite WIP. Ce mécanisme a permis
        d'absorber les trois incidents survenus pendant le projet --- urgence disque,
        AlertManager silencieusement cassé et panne de VictoriaMetrics --- sans casser
        la planification du sprint en cours.

        \paragraph{Définition de terminé.}
        Une carte ne rejoint la colonne \textit{Done} que lorsque les cinq conditions
        suivantes sont simultanément satisfaites :
        \begin{enumerate}[label=\arabic*., leftmargin=*, itemsep=2pt]
            \item la fonctionnalité est déployée et en exécution sur le VPS de production ;
            \item les logs de l'agent restent propres après le passage du nouveau code, sans exception ni avertissement ;
            \item les métriques Prometheus reflètent son comportement lorsque cela est applicable ;
            \item elle a été déclenchée par un événement réel (webhook PR, alerte Prometheus ou requête chat) et a produit le résultat attendu ;
            \item la section correspondante de la documentation est à jour.
        \end{enumerate}
        Cette définition stricte évite que des fonctionnalités à moitié intégrées ne s'accumulent dans le système.

    \subsection{Illustration du tableau Scrumban en pratique}

        L'outil Jira a servi de support à la gestion quotidienne du tableau Scrumban
        tout au long du projet : création des cartes, suivi des transitions de
        colonne, application des limites WIP et association aux jalons. La figure
        \ref{fig:jira_sprint} présente une capture du tableau en fin de Sprint 3,
        retenu ici comme exemple représentatif de l'application concrète de la
        méthodologie. Cette capture donne à voir trois éléments structurants : la
        limite WIP=1 sur la colonne \textit{In Progress}, qui a contraint les cartes
        à être traitées une à une malgré leurs dépendances mutuelles ; les six
        colonnes du flux \textit{Backlog} $\rightarrow$ \textit{Ready}
        $\rightarrow$ \textit{In Progress} $\rightarrow$ \textit{In Review}
        $\rightarrow$ \textit{Testing} $\rightarrow$ \textit{Done} effectivement
        utilisées ; et la matérialisation du jalon M2 par une carte étiquetée
        \textit{Fixed Date}. Le détail technique du Sprint 3 lui-même --- objectif,
        obstacles rencontrés et livrable --- est documenté au chapitre 4.

        \begin{figure}[H]
            \centering
            \frame{\includegraphics[width=\textwidth]{jira_sprint_s3}}
            \caption{Tableau Jira en fin de Sprint 3 --- application concrète
                     des règles Scrumban : WIP=1 sur \textit{In Progress},
                     flux à six colonnes et carte de jalon en classe de
                     service \textit{Fixed Date}.}
            \label{fig:jira_sprint}
        \end{figure}

    \subsection{Déroulement du projet}

        Le projet s'est déroulé sur cinq mois, du 2 février au 30 juin
        2026, structurés en dix sprints de deux semaines : huit sprints
        majeurs (Sprint~1 à Sprint~8, du 2 février au 24 mai) suivis de deux
        micro-sprints d'extension (Post-Sprint~8 et Sprint~8.1, du 25 mai au
        21 juin) consacrés à l'évaluation d'un backend
        d'inférence alternatif, à la robustesse en production et à la
        préparation de la soutenance. La rédaction du rapport a été conduite
        en parallèle des sprints 6 à 10, à raison d'une heure par jour
        bloquée en fin d'après-midi.

        Six jalons rythment cette progression et constituent les points
        d'engagement pris vis-à-vis de l'encadrement. Le jalon M1 marque
        l'existence d'un pipeline vivant, capable d'analyser une
        \textit{Pull Request} réelle de bout en bout. Le jalon M2 certifie
        l'intelligence complète du système avec les cinq scanners et la
        revue LLM publiée sur GitHub. Le jalon M3 valide l'autonomie
        opérationnelle, c'est-à-dire la capacité du système à émettre des
        alertes, à publier sur Slack et à exécuter son scheduler sans
        intervention humaine. Le jalon M4 confirme l'observabilité
        complète. Le jalon M5 atteste le durcissement de production après
        l'audit VPS et la réécriture du parser. Enfin, le jalon M6 valide
        la robustesse de la plateforme sur quatre semaines de
        fonctionnement continu, juste avant la soutenance.

        Le diagramme de Gantt de la figure \ref{fig:gantt} synthétise cette
        progression sur l'axe temporel. Les dix sprints y apparaissent sous
        forme de barres horizontales, les six jalons sous forme de losanges
        positionnés en fin de sprint correspondant, et la rédaction du
        rapport sous forme d'une bande continue chevauchant les derniers
        sprints. La lecture verticale met en évidence les transitions de
        phase, notamment le passage du développement aux micro-sprints
        d'extension entre Sprint~8 et Post-Sprint~8. La lecture horizontale
        révèle quant à elle la cadence régulière des livraisons sur
        l'ensemble de la période.

        \begin{figure}[H]
            \centering
            \frame{\includegraphics[width=\textwidth]{gantt_chart}}
            \caption{Diagramme de Gantt du projet — dix sprints, six jalons
                     et rédaction parallèle du rapport, du 2 février
                     au 30 juin 2026.}
            \label{fig:gantt}
        \end{figure}

\section*{Conclusion}
\addcontentsline{toc}{section}{Conclusion}

Dans ce premier chapitre, nous avons présenté le contexte général du projet au sein de
la Banque de Tunisie et des Emirats, puis nous avons décrit le processus de revue de
sécurité existant et identifié ses principales limites. Cela nous a conduit à présenter
le BTE Security AI Agent comme solution à ces lacunes. Le choix méthodologique a fait
l'objet d'une comparaison entre les principaux frameworks Agile, à l'issue de laquelle
Scrumban a été retenu : sa combinaison de la cadence des sprints héritée de Scrum et
du pilotage par le flux propre à Kanban offre un cadre adapté à un projet solo de cinq
mois en environnement de production. Ce cadre apporte à la fois une visibilité sur les
jalons pour l'encadrement et une souplesse dans la gestion quotidienne lorsque les
incidents arrivent sans prévenir.

Le chapitre suivant porte sur la présentation des différentes technologies sur lesquelles
repose notre solution, avec une étude comparative entre les alternatives disponibles pour
chaque composant, ainsi que la spécification des besoins fonctionnels et non fonctionnels.
